No flat subsidy fits every market. Covering every apartment market at the light benchmark takes \euro{}\attBaseSMaxTwoFiftyApt{} per m\textsuperscript{2}, against a conditional median of \euro{}\attBaseSMedTwoFiftyApt{}, and for houses \euro{}\attBaseSMaxTwoFiftyHouse{} against \euro{}\attBaseSMedTwoFiftyHouse{}; at the medium benchmark the figures are \euro{}\attBaseSMaxFiveHundredApt{} against \euro{}\attBaseSMedFiveHundredApt{} and \euro{}\attBaseSMaxFiveHundredHouse{} against \euro{}\attBaseSMedFiveHundredHouse{}. A uniform national subsidy is therefore poorly targeted because the cost-sharing requirement varies across markets, although $S^*$ is an incentive-compatibility object and not a welfare optimum (Remark~\ref{rem:welfare}). Figure~\ref{fig:policy_map} maps $S^*_m/\hat R$ at the light benchmark, with bubble size encoding the local threshold $\bar R_m$. A flat per-m\textsuperscript{2} grant set at the cross-market mean would over-pay markets where the private incentive already suffices and under-pay those where the gap is largest, the first through fiscal waste and the second through forgone renovations. The dispersion is systematic: the cost-sharing ratio is low where $\bar R_m$ is high, and $\bar R_m$ is high where local prices are high and the energy-efficiency premium $\hat\beta_{1,m}$ is large. The premium reflects the capitalised value of verified quality to sensitive buyers, not a freestanding measure of how many buyers are sensitive (Appendix~\ref{app:rho_identification} shows the two are not identified apart in price data alone). An efficient programme would differentiate the subsidy along these axes and co-target audit-infrastructure investment on markets where the threshold is large but verification is scarce, since $\pi$ does not change the threshold but scales the population it applies to (Section~\ref{ss:att_policy}).
